\section{Proof of Gradient Noise at the Proximal Limit}
\label{app:proximal-gradient-noise}

For Lemma~\ref{lem:proximal-gradient-noise}, fix a prompt $x$, a finite
response space $\mathcal Y_x$, and a positive teacher $\pi_t$ independent
of $\theta$. Assume $\pi_\theta(y\mid x)>0$ is continuously differentiable
on a neighborhood of $\theta_0$, with $\pi_{\theta_0}=\pi_t$.
Write $s_\theta(y\mid x)=\nabla_\theta\log\pi_\theta(y\mid x)$ and
$w_\theta(y\mid x)=\log[\pi_t(y\mid x)/\pi_\theta(y\mid x)]$.
For $m\ge1$ independent samples from the respective distributions, define
\begin{align}
 \widehat g_{\mathrm{KD}}(\theta)
 &=\frac1m\sum_{i=1}^m s_\theta(Y_i\mid x),
 &Y_i&\overset{\mathrm{iid}}{\sim}\pi_t(\cdot\mid x),
 \label{eq:proximal-kd-estimator}\\
 \widehat g_{\mathrm{OPD}}(\theta)
 &=\frac1m\sum_{i=1}^m w_\theta(\widetilde Y_i\mid x)
                          s_\theta(\widetilde Y_i\mid x),
 &\widetilde Y_i&\overset{\mathrm{iid}}{\sim}\pi_\theta(\cdot\mid x).
 \label{eq:proximal-opd-estimator}
\end{align}
The sampled responses and OPD weights are detached during the score update.

\begin{proof}
For every $\theta$, the score identity gives
$\mathbb E_{\pi_\theta}[s_\theta]=\nabla_\theta\sum_y\pi_\theta(y\mid x)=0$.
Hence the per-prompt gradients of the negative-KL objectives are
\begin{align}
 \nabla_\theta[-D_{\mathrm{KL}}(\pi_t\Vert\pi_\theta)]
 &=\mathbb E_{\pi_t}[s_\theta],
 \label{eq:proximal-kd-gradient}\\
 \nabla_\theta[-D_{\mathrm{KL}}(\pi_\theta\Vert\pi_t)]
 &=\mathbb E_{\pi_\theta}[(w_\theta-1)s_\theta]
 =\mathbb E_{\pi_\theta}[w_\theta s_\theta],
 \label{eq:proximal-opd-gradient}
\end{align}
proving unbiasedness. At $\theta_0$, $w_{\theta_0}(y\mid x)=0$ for
every response, so each OPD summand vanishes. Meanwhile,
$\mathbb E_{\pi_t}[s_{\theta_0}]=0$; sample independence yields
\begin{equation}
 \operatorname{Cov}(\widehat g_{\mathrm{KD}}(\theta_0))
 =\frac1m\mathbb E_{\pi_t}[s_{\theta_0}s_{\theta_0}^{\top}]
 =\frac1m F_x(\theta_0).
 \label{eq:app-proximal-fisher}
\end{equation}
Consequently, $\mathbb E\|\widehat g_{\mathrm{KD}}(\theta_0)\|_2^2
=\operatorname{Tr}F_x(\theta_0)/m>0$ whenever the Fisher trace is positive.
This implies positive probability of a nonzero KD estimate, rather than
nonzero estimates for every possible batch.
\end{proof}

\paragraph{A neighborhood of the limit.}
For $A\in\{\mathrm{KD},\mathrm{OPD}\}$, define the estimation error
$V_A(\theta)=\mathbb E\|\widehat g_A(\theta)
-\mathbb E\widehat g_A(\theta)\|_2^2$.
The finite-sum expressions
\begin{align}
 mV_{\mathrm{KD}}(\theta)
 &=\mathbb E_{\pi_t}\|s_\theta\|_2^2
       -\|\mathbb E_{\pi_t}s_\theta\|_2^2,\nonumber\\
 mV_{\mathrm{OPD}}(\theta)
 &=\mathbb E_{\pi_\theta}[w_\theta^2\|s_\theta\|_2^2]
       -\|\mathbb E_{\pi_\theta}[w_\theta s_\theta]\|_2^2
 \label{eq:proximal-estimation-error}
\end{align}
are continuous. Since $V_{\mathrm{OPD}}(\theta_0)=0$ and
$V_{\mathrm{KD}}(\theta_0)=\operatorname{Tr}F_x(\theta_0)/m$,
a positive Fisher trace implies $V_{\mathrm{OPD}}<V_{\mathrm{KD}}$
throughout some neighborhood of $\theta_0$.

\paragraph{Estimator scope.}
The zero-noise conclusion uses the log-ratio estimator in
\eqref{eq:proximal-opd-estimator}: retaining the zero-mean score term
$-s_\theta$ in \eqref{eq:proximal-opd-gradient} gives another unbiased
estimator, but it need not vanish samplewise at the limit. Likewise,
exact teacher-distribution cross-entropy, including full-vocabulary KD at
each sampled prefix, has zero gradient when teacher and student coincide.
The comparison concerns sampled-response KD, with a fixed teacher and
unmodified sampling distributions. Averaging independent prompt-response
pairs over the fixed $\mathcal D$ preserves the limit result, replacing
$F_x$ by $\mathbb E_{x\sim\mathcal D}F_x$ when equality holds for every prompt.
